\documentclass{article}
\usepackage{scrextend}
\usepackage{parskip}
\usepackage{microtype}
\usepackage{geometry}
\usepackage{footmisc}
\usepackage{hyperref}
\usepackage{fancyhdr}
\usepackage{enumitem}

%\geometry{top=1in, bottom=1in, left=1in, right=1in}
%\setlength\parindent{0pt}
%\setlength{\parskip}{6pt}

% Section numbering: Roman numerals at top level, Roman.Alpha at second level
%\renewcommand{\thesection}{\Roman{section}}
%\renewcommand{\thesubsection}{\Roman{section}.\Alph{subsection}}

% Match the spare section style of the reference document
%\usepackage{titlesec}
%\titleformat{\section}{\normalfont\normalsize\bfseries}{\thesection.}{0.5em}{}
%\titlespacing*{\section}{0pt}{14pt}{4pt}
%\titleformat{\subsection}{\normalfont\normalsize\bfseries}{\thesubsection}{0.5em}{}
%\titlespacing*{\subsection}{0pt}{10pt}{3pt}

% Footnotes flush with margin

\setlength{\footnotemargin}{1.5em}

% Page numbers centered at bottom, no header rule
\pagestyle{fancy}
\fancyhf{}
\renewcommand{\headrulewidth}{0pt}
\fancyfoot[C]{\small\thepage}

% Compact lists
\setlist[itemize]{noitemsep, topsep=3pt, leftmargin=1.5em}

% ─────────────────────────────────────────────────────────────────────────────
\begin{document}

% Memo header — plain tabular, matching reference document's spare aesthetic
\begin{tabular}{@{}ll}
\textbf{To:}      & Kandyce Mumm              \\[3pt]
\textbf{From:}    & Jonathan Erdmann          \\[3pt]
\textbf{Subject:} & Proposal: Strategic Career Development Framework Pilot \\[3pt]
\textbf{Date:}    & May 29, 2026                                           \\
\end{tabular}

%\noindent\rule{\textwidth}{0.4pt}
\vspace{4pt}

% ─────────────────────────────────────────────────────────────────────────────
\section{Executive Summary}

Mission Systems conducts an annual individual development planning cycle that
represents a meaningful investment of manager and employee time across the
organization. Based on direct observation through one-on-one conversations and
consistent findings across pulse survey data, the return on that investment is
insufficient. Engineers are producing development plans that are shallow in
quality, limited to a one-year horizon, and rarely executed against in any
substantive way.

This memo proposes a structured pilot program designed to address that gap
across three dimensions: 

\begin{addmargin}[2em]{3em}

\textbf{Plan Quality and Executability}: By providing senior 
	engineers with a rigorous, research-grounded framework for
	building a genuine career strategy rather than a compliance document. 

\textbf{Mentorship and Peer Development Activation}: The organization contains
	significant relational and institutional capital at the senior level that is not
	being converted into productive developmental relationships at the scale our
	depth of talent should enable.

\textbf{Organizational Development Velocity}: Senior engineers who 
	cannot articulate their trajectories develop more
	slowly, are harder to deploy strategically, and represent a 
	compounding cost to succession readiness and retention over time. 

\end{addmargin}

Each dimension is a symptom of the same underlying condition: the 
absence of a structured methodology that gives senior engineers the tools to translate 
their domain expertise into a deliberate, executable career strategy. 

The framework proposed here is grounded in validated research across career
development, goal orientation, and professional identity literature, and has been
designed with the specific context of senior technical professionals in
program-driven environments in mind.  The proposed framework is a structured method 
for producing better outputs from time the organization is already committing. 
The author has applied the framework to his own career development plan, attached in anonymized
form as \hyperref[AppendixC]{Appendix C}, as a demonstration of the
method and the quality of output it produces.

The proposed pilot is structured as follows: a cohort of 10 high-potential senior
engineers (or above), selected in coordination with leadership, organized into five peer
accountability pairs that will work collaboratively to build and
pressure-test individual career strategies over a six-month period. 
Defined success criteria will be assessed at the
conclusion of the pilot, and a findings report delivered to leadership will
include a recommendation on whether and how to scale the approach across the
broader organization.

The proposed next step is a 30-minute conversation to align on cohort selection,
timeline, and organizational dependencies before launch, and to discuss the
value of senior leadership sponsorship in signaling the program's organizational standing
to pilot participants.

% ─────────────────────────────────────────────────────────────────────────────
\section{The Problem}

\subsection{Return on Administrative Overhead}

The annual individual development planning cycle represents one of the more
significant recurring investments of administrative overhead that the department 
makes in its engineering workforce. Managers and engineers commit
meaningful hours each cycle to a process designed in principle to produce
clarity of direction, accelerate capability development, and strengthen the
organization's talent pipeline. In practice, the outputs of that process fall
consistently short of those objectives.

The gap is not attributable to a lack of effort or intent on the part of either
managers or engineers. It is attributable to the absence of a framework capable
of producing the quality of output the process is ostensibly designed to
generate. Without that structure, the annual cycle defaults to what is
immediately legible - near-term skill gaps, current role performance, and
training requests - rather than what is strategically valuable: a genuine,
executable career strategy grounded in a clear understanding of where the
individual is going and what it will take to get there. The result is an
organization that spends the overhead without capturing the return.

\subsection{Plan Quality and Execution Gap}

The pattern this plan addresses is a structural condition and the natural outcome 
of asking technically excellent people to perform a strategic planning exercise 
without the tools to do it well. Engineers who are exceptionally capable of 
solving complex technical problems within defined parameters are being asked to define their own
parameters, articulate their own trajectory, and construct a multi-year strategy
for their professional development without a methodology, a
framework, and without a worked example of what a high-quality output looks
like. The gap in plan quality is a predictable consequence of this condition and
not a reflection of the capability of the individuals experiencing it.

\subsection{Network and Mentorship Underutilization}

Mission Systems carries significant mentorship potential within its senior
engineering and leadership ranks including technical depth, program experience,
and professional networks that represent genuine developmental value for the next generations of
senior talent. This potential is not being converted into productive
relationships at the scale the organization's depth of expertise should enable.

Mentorship relationships form most naturally when both parties share a language for articulating
development goals, identifying relevant experience, and defining what a
productive relationship would produce. Without a common framework, the
conditions for those conversations to begin are absent. Potential mentors have
no clear mechanism for engaging, and potential mentees have not yet developed
the strategic clarity to make a specific, meaningful ask. The relationships that
do form tend to be informal, uneven in quality, and concentrated among engineers
who already possess the professional confidence to seek them out, precisely
the engineers who need them least.

\subsection{Organizational Development Velocity}

The consequences of the gap described above are compounding in nature.
An organization whose senior engineers cannot articulate their trajectories
develops more slowly than one that can, not because the individuals are less
capable, but because deliberate development is more efficient than incidental
development. Engineers who understand where they are going make better decisions
about where to invest their discretionary effort, which stretch assignments to
seek, which relationships to build, and which capabilities to develop ahead of
demand rather than in response to it.

At the senior engineer level, where the organization's investment per
individual is highest and the lead time for developing strategic capability is
longest, the cost of that inefficiency is most consequential. The
three-to-five year horizon is where the compounding effect of deliberate versus
incidental development becomes visible in succession readiness, capability
depth, and retention. A pilot program that begins generating that compounding
effect now produces organizational returns that a program launched two years
from now cannot fully recover.

% ─────────────────────────────────────────────────────────────────────────────
\section{The Proposed Framework}

The framework is a three-layer strategic career management system that converts
the existing IDP cycle from a compliance exercise into an executable personal
strategy. The framework is designed specifically for senior technical professionals in
program-driven environments and not adapted from a generic talent development
instrument, but constructed from the ground up with the specific challenges of
that population in mind. The terminal output is a Strategic Career Plan that is
qualitatively different from our standard IDP in specificity, time horizon, and
operational rigor. The full framework definition is provided in 
\hyperref[AppendixA]{Appendix A}.

\subsection{Framework Layers}

The framework operates across three layers with distinct purposes and time horizons.
\begin{addmargin}[2em]{3em}

\textbf{Strategic Foundation}: Defines where the engineer is going and why they 
	will thrive there. The strategy statements contain the enduring elements that anchor 
	everything else within the plan: personal mission, positioning thesis, target functions, 
	known gaps, and breakthrough goals. 

\textbf{Operating Model}: Translates that strategy into an executable, milestone-gated plan organized 
	around annual objectives, initiatives, and action items. Each element traces directly to a parent,
	ensuring nothing in the plan exists without strategic justification. 

\textbf{Performance Cadence}: Drives accountability and surfaces signals through four structured weekly,
	monthly, quarterly, and annual review loops. Each review asks a distinct question about execution
	and strategy alignment.

\end{addmargin}

\newpage

\subsection{Course Correction Mechanism}

The framework includes an explicit decision rule that distinguishes between problem types:
\begin{addmargin}[2em]{3em}
	\textbf{Execution Problems}: Actions are \textit{not} being completed and the Operating Model requires updating.

	\textbf{Strategy Problems}: Actions \textit{are} being completed but leading indicators are not moving and the
	Strategic Foundation requires recalibration.
\end{addmargin}

This mechanism prevents the two most common failure modes in personal development planning: the rigid plan that 
never adapts, and the reactive plan that never generates signal.

\subsection{Research and Methodological Foundation}

The framework draws on validated research across career development,
professional identity, and goal orientation literature, and adapts its
operational architecture from Hoshin Kanri strategy deployment, phase-gate
program management, and Work Breakdown Structure principles. The full
theoretical foundation is summarized in \hyperref[AppendixB]{Appendix B}.

%\footnote{The application of
%this operational architecture to individual career strategy --- including the
%correction mechanism and traceability requirement --- represents an original
%adaptation not present in existing personal career development literature.}

% ─────────────────────────────────────────────────────────────────────────────
\section{Pilot Design}

\subsection{Scope}

The pilot is proposed as a six-month program engaging a cohort of 10
high-potential senior engineers selected in coordination with leadership. Cohort
selection will prioritize engineers who represent the range of tenure, technical
discipline, and career stage present across the broader deparment engineering population 
and not exclusively the most advanced, but a cross-section capable of generating findings 
relevant to a broader rollout. The cohort will be organized into five peer accountability 
pairs for the duration of the pilot, with pairing designed to maximize complementary 
experience profiles and developmental value for both participants.

The pilot will be anchored to the existing annual IDP cycle wherever timing
permits, with the intent of replacing rather than supplementing the standard IDP
conversations for pilot participants. The objective is to demonstrate that the
framework produces a superior output within the same time commitment the
organization is already making.

\subsection{Implementation Approach}

The pilot will be delivered through a series of working sessions in which
participants build their career strategies in real time rather than receiving
instruction and completing work independently afterward. Between sessions, each
peer pair will be responsible for continuing to develop and refine their
individual plans, using their partner as a sounding board and a source of honest
challenge. The facilitator will conduct structured check-ins at the two-month
and four-month marks to assess plan development and make any necessary
adjustments before the concluding session. Full implementation detail is
provided in \hyperref[AppendixD]{Appendix D}.

\subsection{Success Criteria}

The pilot will be evaluated against four criteria at the six-month mark.


\begin{addmargin}[2em]{3em}
\textbf{Plan Quality}: Strategic Career Plans will be assessed for specificity,
time horizon depth, and actionability relative to IDP outputs those same
engineers produced in prior cycles representing a within-individual comparison that
controls for variation in starting capability across the cohort.

\textbf{Execution Rate}: Facilitator assessment at two-month and
four-month check-ins will document whether participants are actively working
against their plans rather than treating the completed document as a terminal
artifact.

\textbf{Relationship Formation}: New mentorship and peer development
relationships formed among and beyond the cohort during the six-month period,
as a direct measure of the framework's effectiveness in activating the
relational potential the organization currently underutilizes.

\textbf{Participant Assessment}: Structured qualitative feedback from each
participant at pilot conclusion, informing the findings report and any
recommendation to scale.
\end{addmargin}


\subsection{Reporting Commitment}

At the conclusion of the six-month pilot, a findings report will be delivered
to the sponsoring leader documenting performance against each success criterion,
aggregating participant feedback, and making a specific recommendation on
whether and how to scale the program across the broader 
engineering organization. The findings report will not advocate for a
predetermined conclusion. If the pilot surfaces evidence that the framework
does not produce the outcomes proposed here, that finding will be reported
accurately and a recommendation against formally implementing the plan will be
provided.

% ─────────────────────────────────────────────────────────────────────────────
\section{Request}

The author requests authorization to conduct a six-month pilot of the Strategic
Career Development Framework as scoped in the above sections, with a cohort of 10
high-potential senior engineers selected in coordination with sponsoring leadership,
organized into five peer accountability pairs, and anchored to the existing
annual IDP conversations.

Two elements would meaningfully strengthen the pilot's organizational standing
and increase the quality of its findings. The first is senior leadership-level sponsorship 
to serve as a visible endorsement that signals to participants that the program 
carries organizational weight and that the plans they produce will be seen and 
considered by leadership. The second is alignment on cohort selection criteria before launch,
to ensure the cross-section of participants is representative enough to generate findings 
relevant to a broader rollout.

The framework has been applied to the author's own career development plan,
attached as \hyperref[AppendixC]{Appendix C}, as a demonstration of the methodology and the quality
of output it produces.

A 30-minute conversation to align on cohort selection, timing, and sponsorship
before launch is proposed as the next step.

\newpage
\appendix

\section{Framework Definition}
\label{AppendixA}

\newpage
\section{Research Foundation}
\label{AppendixB}

\newpage
\section{Author's Anonymized Strategic Career Development Plan}
\label{AppendixC}
%\documentclass{article}
\usepackage{amsmath}
\usepackage{mathtools}
\usepackage{amsfonts}
\usepackage{parskip}
\setlength\parindent{0pt}

\title{Career Development Plan: Layer 1 - Strategic Foundation - DRAFT \\ [1em]
	\large Jonathan Erdmann \\ [1em]
}

\begin{document}

\maketitle

This document is the strategic foundation of my Career Development Plan and serves as the anchor for all planning activity that 
follows. Content in this layer operates at the level of enduring truth.
Statements here should remain stable across the full planning horizon and require deliberate annual review before they are changed.
Every element in Layers 2 and 3 must be traceable to something written here.

\section{Personal Mission}

To apply deep domain expertise and rigorous analytical thinking to complex, high-stakes problems in aerospace and defense, 
owning the full arc from problem definition to decision, such that consequential decisions can be directly traced to the 
quality of my work.

\section{Positioning Thesis}

Two decades in aerospace and defense across structural integrity programs, mission-critical systems, engineering organizations,
and workforce management at scale has produced deep domain fluency, technical judgment, and a firsthand understanding of how
programs get executed and how organizations perform. An Executive MBA in finance and strategy adds to the quantitative and analytical
toolkit to operate in corporate strategy, corporate development, and investment functions. Together these form a professional profile
that can move fluidly between technical substance, organizational dynamics, and strategic implications. These capabilities 
are a unique and valuable combination and are consistently relevant in aerospace and defense.

\section{Target Functions}

%Target functions are sequenced deliberately. Near-term targets inform the operating model immediately. Medium and long-term targets
%inform capability development decisions and are actively revisited at each annual review.

\subsection{Near-Term Primary Target: Corporate Strategy}

The internal advisory function supporting C-suite decision-making at large firms. Owns strategic planning cycles, competitive
intelligence, market sizing, portfolio tradeoffs, and build/buy/partner analysis. Output is generally recommendations and 
frameworks, not transactions or operational decisions. The work is analytical and influence-driven, shaping decisions without 
always owning their execution. At a large prime this function sits close to the CEO and business unit presidents, making 
it one of the highest-visibility staff roles available to someone without a P\&L.

\textbf{Entry Angle:} The positioning thesis directly supports this function --- domain fluency, technical judgment, and quantitative
finance training are precisely differentiators of a strong corporate strategy candidate at a defense prime from a 
generalist consultant. The gap to close is the strategy title itself and demonstrated exposure to strategic planning processes, 
which the operating model must address explicitly.

\subsection{Medium-Term Target: Corporate Development}

This is the transaction-execution function: Owns M\&A sourcing, financial modeling, due diligence, deal structuring, 
and integration planning. Success is measured in closed transactions and not recommendations to decision makers.
At a large prime this function operates with significant financial authority and works directly with executive 
leadership and external advisors on inorganic growth decisions.

\textbf{Entry Angle:} Direct entry without prior corporate development or investment banking experience is uncommon. 
The credible path is through corporate strategy first, where exposure to build/buy/partner analysis, 
portfolio reviews, and strategic transactions creates a natural bridge. The operating model should include 
deliberate M\&A adjacent activities --- whether internal transaction support, partnership structuring, 
or strategic initiative ownership --- that begin building this credential during Phase 2.

\newpage
\subsection{Long-Term Option: Corporate Venture Capital} 

The investment function focused on minority stakes in early-to-growth stage companies, typically housed within 
a prime's dedicated investment arm. Balances strategic fit activities like technology scouting, capability 
adjacencies, ecosystem development with financial return. In aerospace and defense, CVC often functions 
as a forward-deployed R\&D pipeline, identifying technologies before they are mature enough for acquisition. 
Investment theses are evaluated on both strategic relevance and commercial viability.

\textbf{Entry Angle:} Technical domain depth in aerospace and defense systems is a genuine differentiator in this function --- the
ability to evaluate whether a technology is real, a team is credible, and a capability is strategically relevant may be
rare among people with the business credentials to operate in an investment role. The credible path runs through 
corporate development, where deal evaluation and financial modeling skills are formalized, or through sufficient 
ecosystem visibility that a CVC role becomes a direct lateral. 
Preserved as optionality and not actively targeted until Phase 2 is complete.

\newpage

\section{Known Gaps}

Gaps are assessed across two domains --- competency and positioning --- and three time horizons corresponding 
to the target function sequence. Near-term gaps are assessed in detail. Medium and long-term gaps are 
assessed directionally and revisited at each annual review.

Competency gaps address the ability to perform and thrive once in the role. Positioning gaps address the ability to be considered for
and selected into the role. Both must be addressed. Closing one at the expense of the other produces either an unprepared candidate
or an invisible one.

Competency is assessed first because execution of the mission depends on the quality of the work, not the quality of the positioning.

\subsection{Near Term: Corporate Strategy}

\subsubsection{Competency Gaps}

\textbf{Strategic Analysis \& Frameworks}
Some exposure through coursework and work experience but not yet practiced independently at the level required to structure and solve
ambiguous, high-stakes problems without guidance. Formal frameworks have been introduced through the EMBA but applied experience in a
corporate strategy context is non-existent. This is the most critical competency gap to close given the centrality of 
structured analytical thinking to the target function.

\textbf{Financial Acumen}
Foundational understanding established through EMBA coursework but limited applied experience building, stress-testing, or presenting
financial models in a business context. Fluency with financial statements and capital allocation logic exists at a conceptual level but
has not yet been exercised in a corporate strategy or corporate development setting.

\textbf{Industry \& Domain Knowledge}
Deep operational knowledge from within the industry but limited breadth across the competitive landscape, program economics at the
enterprise level, and technology trends outside areas of direct experience. The domain depth is a genuine asset; the gap is in breadth
and in the ability to synthesize across the landscape rather than within a specific program or system.

\textbf{Executive Communication}
Capable of structuring and delivering findings to senior leadership with significant room to grow at more senior levels. 
The gap is in consistency, precision under scrutiny, and the ability to compress complex analysis into the format and 
cadence that corporate strategy demands. Significant gap when communicating in settings where dominant voices frequently
insert themselves.

\textbf{Stakeholder Navigation}
Capable of navigating organizational complexity and building credibility across functions without direct authority. The gap is in
operating at the corporate level where political complexity, ambiguity of authority, and the stakes of decisions are 
materially higher.

\subsubsection{Positioning Gaps}

\textbf{Professional Narrative}
A narrative exists grounded in two decades of A\&D experience but is not yet coherent or tailored to corporate strategy. Current
profile and resume position as an engineering manager, not a strategist. The work required is reframing existing experience ---
which includes genuine strategic work products --- into a story that is legible and compelling to a decision-maker in the target function.

\textbf{External Visibility}
A LinkedIn profile exists but is not yet visible to the right people in the target function. No presence in corporate strategy or
corporate development circles at large primes. The gap is both reach and relevance.

\textbf{Network}
Broad professional network but not targeted to corporate strategy or corporate development functions at large primes. Limited
relationships with people who could surface my name when an opportunity arises or provide insider context on how these functions operate.

\textbf{Internal Positioning}
Some visibility to senior leadership at Collins through current role performance but not yet associated with strategic capability or
ambition beyond engineering management. No executive sponsor actively advocating for development or advancement into a strategy-adjacent role.

\textbf{Credentials}
EMBA in progress; completion expected May 2027. No strategy-adjacent title or role held to date. Like professional 
narrative, part of this gap is due to an engineering-first background.

\subsection{Medium-Term: Corporate Development}

To be assessed directionally following completion of near-term gap closure. Key anticipated gaps include transaction experience,
financial modeling depth, and M\&A process knowledge. Full assessment to be conducted at the Phase 2 annual review.

\subsection{Long-Term: Corporate Venture Capital}

To be assessed at a high level following completion of medium-term gap closure. Key anticipated gaps include investment thesis
development, startup evaluation, and portfolio management experience. Full assessment to be conducted at the Phase 3 annual review.

\newpage

\section{Breakthrough Goals}

\textbf{Goal 1: Complete the SMU Executive MBA}

This is the most time-bound and definitive goal in the plan. Completion closes the credential gap and activates the
quantitative finance and strategic analysis toolkit required to perform in the target function. Graduating with distinction 
signals the seriousness of the credential to a hiring decision-maker although is not necessary and is harder given difficulties
in Term 3 due to Achilles tendon rupture.

\textbf{Goal 2: Establish a coherent, credible professional narrative tailored to corporate strategy and a targeted network
of relevant relationships} 

Success condition: the right people inside Collins and within the broader A\&D corporate strategy community know who I am, understand
what I am building toward, and would surface my name when a relevant opportunity arises. Internal positioning takes priority in the
near term given the long-term Collins-first strategy.

\textbf{Goal 3: Transition into a strategy-adjacent role or formal corporate strategy function within Collins Aerospace
at the Senior Manager level or above}

Collins is the first-move target because existing relationships, organizational knowledge, and demonstrated performance reduce the
credibility gap that an external employer would require evidence to close. An external move to a corporate strategy function at
another company within RTX is the contingency if internal advancement stalls, assessed at the annual review.

\textbf{Goal 4: Own at least one consequential strategic work product within Collins Aerospace where the output directly
informed a senior executive decision.}

Success condition: a market assessment, portfolio recommendation, competitive analysis, or strategic initiative with a traceable
outcome where the link between the work and the decision is direct and attributable.

\textbf{Goal 5: Be positioned credibly for a corporate development role at a large A\&D prime, with M\&A adjacent exposure
and financial modeling depth sufficient to compete as a serious candidate.}

This goal preserves the medium-term target without conflating it with the near-term transition. 
The goal is achieved through the combination of corporate strategy credentialing, deliberate pursuit of M\&A adjacent 
activities, and deepening financial acumen during Phases 2 and 3.

\end{document}


\newpage
\section{Pilot Program Implementation Plan}
\label{AppendixD}

% 
%\vspace{20pt}
%\noindent\rule{\textwidth}{0.4pt}
%\vspace{4pt}

% Appendix placeholders
%\renewcommand{\thesection}{\Alph{section}}
%\setcounter{section}{0}
%\titleformat{\section}{\normalfont\normalsize\bfseries}{Appendix \thesection.}{0.5em}{}
%
%\section{Framework Definition}
%
%\textit{The complete three-layer framework definition --- governing principle,
%Strategic Foundation, Operating Model, and Performance Cadence --- with
%component definitions, element hierarchy, leveling guidance, and correction
%mechanism. To be attached as a separate document.}
%
%\section{Research Foundation}
%
%\textit{Summary of the theoretical pillars underpinning the framework: career
%capital theory, professional identity theory, goal orientation research, and
%the operational architecture sources --- Hoshin Kanri, phase-gate program
%management, and Work Breakdown Structure methodology. Key citations and brief
%annotations for each. To be attached.}
%
%\section{Author's Strategic Career Development Plan}
%
%\textit{The author's Strategic Career Development Plan, produced using the
%framework described in Section~III. Provided as a worked example of the
%methodology and the quality of output it generates. To be attached.}
%
%\section{Pilot Design Detail}
%
%\textit{Full working session structure, facilitation approach, peer pairing
%logic, check-in cadence, and success criteria definitions. To be attached.}

\end{document}
